%======================================================================
\section{Algoritmos}\label{sec:algoritmos}
%======================================================================

\subsection{Arquitetura}

Os três algoritmos são genéricos e interagem com o problema apenas por uma interface comum (Tabela~\ref{tab:interface}). Essa separação permite executar as 9 combinações (3 algoritmos $\times$ 3 instâncias) sem alterar o código dos algoritmos e centralizar a contagem de avaliações em um único ponto: cada chamada a $h$ incrementa um contador único, inclusive a reavaliação da solução corrente.

\begin{table}[H]
\centering
\caption{Interface entre algoritmos e problema.}
\label{tab:interface}
\small
\begin{tabular}{p{3.2cm}p{4.6cm}}
\toprule
Método & Papel\\
\midrule
\texttt{solucao\_aleatoria()} & gera uma permutação inicial\\
\texttt{h(s)} & fitness (minimização)\\
\texttt{sucessores(s)} & itera sobre todos os sucessores\\
\texttt{sucessor\_aleatorio(s)} & sorteia um sucessor\\
\texttt{cruzamento(s1,s2)} & recombinação\\
\texttt{mutacao(s)} & mutação\\
\bottomrule
\end{tabular}
\end{table}

\pend{linguagem/versão, bibliotecas, gerador aleatório e política de sementes (uma semente por execução, registrada em disco)}

\subsection{Avaliação incremental}\label{sec:incremental}

Avaliar $h$ do zero custa $O(n)$. Ao aplicar um movimento cuja menor posição afetada é $i$, o prefixo $(\pi_1,\dots,\pi_{i-1})$ permanece inalterado, de modo que $C_{\pi_{i-1}}$ e o custo acumulado até $i-1$ podem ser reaproveitados a partir de vetores de prefixo mantidos para a solução corrente. Basta recalcular o sufixo, a um custo $O(n-i+1)$. Isso vale para \textit{swap}, \textit{insertion} e 2-opt: mesmo que os travessias sejam simétricas, os instantes $C_j$ dentro do trecho alterado mudam e precisam ser recalculados. A avaliação incremental não altera a contagem de avaliações, apenas o tempo de parede.

\pend{vizinhança(s) com avaliação incremental implementada, descrição do teste automático de concordância com a avaliação completa e seu resultado}

O ganho de tempo medido é reportado na Seção~\ref{sec:res-incremental}.

\subsection{(A) \textit{Greedy descent} com \textit{random restart}}\label{sec:greedy}

A partir de uma permutação aleatória, o algoritmo move-se repetidamente para um sucessor que melhora $h$, até atingir um \textbf{ótimo local}: uma permutação sem sucessor melhor na vizinhança adotada. Ao atingir o ótimo local, aplica-se \textit{random restart}: recomeça-se de nova permutação aleatória, preservando a melhor solução já vista, até esgotar o orçamento $B$. As variantes de seleção do sucessor são:

\begin{itemize}
  \item \textbf{\textit{Best improvement}:} varre toda a vizinhança e escolhe o melhor sucessor;
  \item \textbf{\textit{First improvement}:} aceita o primeiro sucessor que melhora $h$;
  \item \textbf{Sucessor aleatório com aceitação de melhora:} sorteia um único sucessor e o aceita somente se melhorar $h$.
\end{itemize}

\pend{variantes efetivamente implementadas (mínimo duas) e a variante adotada nos experimentos principais, com justificativa}

Sob orçamento fixo, as variantes diferem fortemente em passos de melhoria por avaliação: o \textit{best improvement} paga uma vizinhança inteira por passo, enquanto o \textit{first improvement} e a variante aleatória pagam em média menos. A Seção~\ref{sec:res-t2} reporta quantos passos de melhoria cada variante deu dentro do orçamento. Para a variante aleatória, o critério de ótimo local não é determinístico; é preciso definir o critério de reinício.

\pend{critério de detecção de ótimo local/reinício da variante aleatória, se implementada}

\subsection{(B) \textit{Simulated annealing}}\label{sec:sa}

Seja $\pi$ a solução corrente e $\pi'$ um sucessor sorteado ao acaso. Se $h(\pi')\le h(\pi)$, aceita-se $\pi'$; caso contrário, aceita-se com probabilidade
\begin{equation}
  P=\exp\!\left(-\frac{h(\pi')-h(\pi)}{T}\right),
  \label{eq:gibbs}
\end{equation}
a distribuição de Gibbs. Cada sucessor examinado consome uma avaliação.

\textbf{\textit{Annealing schedule}.} Adota-se resfriamento geométrico $T\leftarrow\alpha T$. Como cada iteração consome uma avaliação e o orçamento é $B$, impor $T_f=T_0\,\alpha^{\,B/m}$ com atualização a cada $m$ iterações fornece
\begin{equation}
  \alpha=\left(\frac{T_f}{T_0}\right)^{m/B}.
  \label{eq:alpha}
\end{equation}
Assim, a temperatura chega a $T_f$ exatamente ao esgotar o orçamento, e o resfriamento é ajustado a $B$ de cada instância.

\textbf{Calibração de $T_0$ e $T_f$.} A escala de $T$ depende da magnitude típica de $h(\pi')-h(\pi)$, que varia de dezenas a milhares e muda com a instância; portanto, os parâmetros são calibrados \textit{por instância}. O procedimento adotado deve ser:

\begin{enumerate}
  \item amostrar um conjunto de pares $(\pi,\pi')$ aleatórios por instância e estimar a distribuição de $\Delta=h(\pi')-h(\pi)$ para $\Delta>0$;
  \item escolher $T_0$ tal que a taxa de aceitação inicial de movimentos piores seja próxima de um alvo alto (quase \textit{random walk});
  \item escolher $T_f$ tal que a taxa de aceitação final seja praticamente nula (quase \textit{greedy descent});
  \item verificar a escolha pelo gráfico da taxa de aceitação ao longo da execução.
\end{enumerate}

\pend{alvos de aceitação, tamanho da amostra de calibração e valores finais de $T_0$, $T_f$, $\alpha$ e $m$ por instância (Tabela~\ref{tab:par-sa})}

\begin{table}[H]
\centering
\caption{Parâmetros do \textit{simulated annealing} por instância.}
\label{tab:par-sa}
\small
\begin{tabular}{lcccc}
\toprule
Instância & $T_0$ & $T_f$ & $\alpha$ & Vizinhança\\
\midrule
pequena       & \pp & \pp & \pp & \pp\\
intermediária & \pp & \pp & \pp & \pp\\
grande        & \pp & \pp & \pp & \pp\\
\bottomrule
\end{tabular}
\end{table}

\textbf{Taxa de aceitação.} Registra-se, em janelas de avaliações, a fração de movimentos \textit{piores} aceitos. Uma têmpera bem calibrada começa quase como \textit{random walk} e termina quase como \textit{greedy descent}; um gráfico que permanece alto indica $T$ alta demais, e um que cai a zero logo no início indica $T$ baixa demais.

\subsection{(C) Algoritmo genético}\label{sec:ga}

O GA mantém uma população de $N$ cromossomos e repete: seleção, cruzamento com probabilidade $P_c$, mutação com probabilidade $P_m$ e elitismo, até esgotar o orçamento de avaliações. A população inicial é uma amostra aleatória de $N$ permutações.

\textbf{Seleção e conversão custo$\to$aptidão.} Como $h$ é de minimização, o operador de seleção deve favorecer valores \textit{menores}. No torneio de tamanho $k$, basta escolher o indivíduo de menor $h$ entre $k$ sorteados, sem conversão. Na roleta, é necessária uma aptidão $f(\pi)\ge 0$ crescente com a qualidade, por exemplo $f(\pi)=h_{\max}-h(\pi)$ (com $h_{\max}$ o pior valor da geração) ou $f(\pi)=1/(1+h(\pi))$.

\pend{operador de seleção adotado, tamanho do torneio ou fórmula de conversão custo$\to$aptidão efetivamente usada, e justificativa}

\textbf{Cruzamento.} Cruzamentos clássicos de um ou múltiplos pontos duplicariam ativos e gerariam cromossomos inválidos; é necessário um operador que preserve a permutação (OX, PMX ou cruzamento por arestas).

\pend{operador de cruzamento adotado (OX, PMX ou arestas) e justificativa}

A título de ilustração, o \textit{Order Crossover} (OX) opera como segue, com $n=8$, pais $P_1=(1\,2\,3\,4\,5\,6\,7\,8)$ e $P_2=(3\,7\,5\,1\,6\,8\,2\,4)$ e pontos de corte antes da posição 4 e depois da posição 6:

\begin{enumerate}
  \item copia-se o segmento de $P_1$ entre os cortes para o filho: $(\_\,\_\,\_\,4\,5\,6\,\_\,\_)$;
  \item percorre-se $P_2$ a partir da posição seguinte ao segundo corte (posição 7), de forma circular: $2,4,3,7,5,1,6,8$;
  \item descartam-se os ativos já presentes no segmento ($4,5,6$), restando $2,3,7,1,8$;
  \item preenchem-se as posições livres a partir da posição 7, também circularmente: posições $7,8,1,2,3$ recebem $2,3,7,1,8$;
  \item o filho é $(7\,1\,8\,4\,5\,6\,2\,3)$, que é uma permutação válida.
\end{enumerate}

\textbf{Mutação.} Aplica-se uma pequena taxa de mutação para inserir diversidade e favorecer a exploração. O ponto de partida é a \textit{swap mutation}; os demais operadores da vizinhança também são admissíveis.

\pend{operador de mutação adotado}

\textbf{Elitismo e parâmetros.} O elitismo preserva os $e$ melhores indivíduos de cada geração. A Tabela~\ref{tab:par-ga} reúne os parâmetros. Cada indivíduo novo avaliado consome uma avaliação, e o algoritmo para exatamente em $B$.

\begin{table}[H]
\centering
\caption{Parâmetros do algoritmo genético.}
\label{tab:par-ga}
\small
\begin{tabular}{lccc}
\toprule
Parâmetro & pequena & interm. & grande\\
\midrule
$N$ (população) & \pp & \pp & \pp\\
$P_c$ & \pp & \pp & \pp\\
$P_m$ & \pp & \pp & \pp\\
Elitismo ($e$) & \pp & \pp & \pp\\
Tamanho do torneio & \pp & \pp & \pp\\
\bottomrule
\end{tabular}
\end{table}

\pend{como os parâmetros foram escolhidos e testados}

\textbf{Diversidade.} A cada geração registra-se o número de cromossomos distintos e a distância média entre pares de indivíduos, definida pela distância de Hamming entre posições $d(\pi,\sigma)=|\{k:\pi_k\neq\sigma_k\}|$. A queda dessas medidas diagnostica a convergência prematura.